% Core mathematical notation.
\def\Gam{\ensuremath{\Gamma}\xspace}
\def\clockSet{\ensuremath{\mathcal{C}}\xspace}
\def\clockValuations#1{\ensuremath{\mathbb{R}_{\ge 0}^{#1}}\xspace}
\def\Nat{\ensuremath{\mathbb{N}}\xspace}
\def\Int{\ensuremath{\mathbb{Z}}\xspace}
\def\Reals{\ensuremath{\mathbb{R}}\xspace}
\def\kI{\ensuremath{I}\xspace}
\def\ck{\ensuremath{c}\xspace}
\def\DBM#1{\ensuremath{\mathbf{D}_{#1}}\xspace}
\def\UP{\ensuremath{\uparrow}\xspace}
\def\RESET#1{\ensuremath{[#1 \mathrel{:=} 0]}\xspace}
\def\tEU#1#2#3{\ensuremath{\mathrm{E}(#1\mathbf{U}_{#3}#2)}\xspace}
\def\tAU#1#2#3{\ensuremath{\mathrm{A}(#1\mathbf{U}_{#3}#2)}\xspace}
\def\EU#1#2{\ensuremath{\mathrm{E}(#1\mathbf{U}#2)}\xspace}
\def\AU#1#2{\ensuremath{\mathrm{A}(#1\mathbf{U}#2)}\xspace}
\def\EX#1{\ensuremath{\mathit{EX}#1}\xspace}
\def\AX#1{\ensuremath{\mathit{AX}#1}\xspace}
\def\tAR#1#2#3{\ensuremath{\mathrm{A}(#1\mathbf{R}_{#3}#2)}\xspace}
\def\tER#1#2#3{\ensuremath{\mathrm{E}(#1\mathbf{R}_{#3}#2)}\xspace}
\def\ER#1#2{\ensuremath{\mathrm{E}(#1\mathbf{R}#2)}\xspace}
\def\AR#1#2{\ensuremath{\mathrm{A}(#1\mathbf{R}#2)}\xspace}
\def\Z{\ensuremath{Z}\xspace}
\def\Deadline{\ensuremath{\mathcal{D}}\xspace}

% Rule-name helpers and clickable cross-references.
\makeatletter
\newcommand{\DeclareRule}[2]{\expandafter\def\csname rule@#1\endcsname{#2}}
\newcommand{\RuleLabel}[1]{{\ensuremath{#1}}}
\newcommand{\UunfRuleName}{\RuleLabel{\mathbf{U}^\bullet}}
\newcommand{\UstopRuleName}{\RuleLabel{\mathbf{U}\text{-stop}}}
\newcommand{\UgoRuleName}{\RuleLabel{\mathbf{U}\text{-go}}}
\newcommand{\UgoCarryRuleName}{\RuleLabel{\mathbf{U}^\ast\text{-go}}}
\newcommand{\RsplitRuleName}{\RuleLabel{\mathbf{R}^{\mathsf{split}}}}
\newcommand{\RbranchName}[1]{\RuleLabel{\mathbf{R}^\bullet_{#1}}}
\newcommand{\RbranchStopName}[1]{\RuleLabel{\mathbf{R}^\bullet_{#1}\text{-stop}}}
\newcommand{\RbranchGoName}[1]{\RuleLabel{\mathbf{R}^\bullet_{#1}\text{-go}}}
\newcommand{\rulename}[1]{\ensuremath{(\csname rule@#1\endcsname)}}
\newcommand{\ruleanchor}[1]{\hypertarget{rule:#1}{\rulename{#1}}}
\newcommand{\ruleref}[1]{\ifcsname r@tag:#1\endcsname\hyperref[tag:#1]{\rulename{#1}}\else\hyperlink{rule:#1}{\rulename{#1}}\fi}
\newcommand{\rname}[1]{\csname rule@#1\endcsname}
\newcommand{\rlink}[1]{\ifcsname r@tag:#1\endcsname\hyperref[tag:#1]{\rulename{#1}}\else\hyperlink{rule:#1}{\rulename{#1}}\fi}
\makeatother
\DeclareRule{AND}{\land}
\DeclareRule{OR}{\lor}
\DeclareRule{AX}{\mathrm{AX}}
\DeclareRule{EX}{\mathrm{EX}}
\DeclareRule{Delay}{\mathrm{Delay}}
\DeclareRule{Conj}{\land}
\DeclareRule{Disj}{\lor}
\DeclareRule{Cov}{\mathrm{Cov}}
\DeclareRule{Covstar}{\mathrm{Cov}^{*}}
\DeclareRule{Uunf}{\UunfRuleName}
\DeclareRule{Runf}{\mathbf{R}^\bullet}
\DeclareRule{Rsplit}{\RsplitRuleName}
\DeclareRule{Runf_before}{\RbranchName{\prec\kI}}
\DeclareRule{Runf_in}{\RbranchName{\in\kI}}
\DeclareRule{Runf_after}{\RbranchName{\succ\kI}}
\DeclareRule{Runf_lapse}{\RuleLabel{\mathbf{R}^\bullet_{\succ\kI}^{\mathsf{lapse}}}}
\DeclareRule{Ustop}{\UstopRuleName}
\DeclareRule{Ugo}{\UgoRuleName}
\DeclareRule{Ugocarry}{\UgoCarryRuleName}
\DeclareRule{Rbeforestop}{\RbranchStopName{\prec\kI}}
\DeclareRule{Rbeforego}{\RbranchGoName{\prec\kI}}
\DeclareRule{Rinstop}{\RbranchStopName{\in\kI}}
\DeclareRule{Ringo}{\RbranchGoName{\in\kI}}
\DeclareRule{Rstop}{\RuleLabel{\mathbf{R}^\bullet\text{-stop}}}
\DeclareRule{Rgo}{\RuleLabel{\mathbf{R}^\bullet\text{-go}}}
\DeclareRule{discharge}{\mathsf{discharge}}
\DeclareRule{launch}{\mathsf{launch}}
\DeclareRule{Rsplitbefore}{\RbranchName{\prec\kI}}
\DeclareRule{Rsplitin}{\RbranchName{\in\kI}}
\DeclareRule{Rsplitafter}{\RbranchName{\succ\kI}}
\DeclareRule{openseg}{\mathsf{open seg.}}
\DeclareRule{boundary}{\mathsf{boundary}}




% Figure helpers.
\newcommand{\myvdots}[4]{%
  % #1 = x-coordinate
  % #2 = y-coordinate of top dot
  % #3 = vertical spacing
  % #4 = dot radius
  \foreach \i [evaluate=\i as \yy using #2-\i*#3] in {0,1,2}
    \fill (#1,\yy) circle[radius=#4];
}

\newcommand{\lemitem}[2]{\hyperref[#1]{\ref*{#1}(\textup{#2})}}

% Shared shorthand commands live in acronyms.tex.
% Acronyms and shorthand commands used across the manuscript.
% Keeping these separate makes the main preamble easier to scan and maintain.

% Semantic shorthands used in the prose.
\providecommand{\eventualities}{\ensuremath{\mathcal{E}}\xspace}
\providecommand{\obligations}{\ensuremath{\mathcal{O}}\xspace}
\providecommand{\ctx}{\ensuremath{\mathit{ctx}}\xspace}
\providecommand{\INF}{\ensuremath{\infty}\xspace}
\providecommand{\SAT}{\ensuremath{\mathrm{SAT}}\xspace}
\providecommand{\UNSAT}{\ensuremath{\mathrm{UNSAT}}\xspace}
\providecommand{\EF}{\ensuremath{\mathrm{EF}}\xspace}
\providecommand{\AF}{\ensuremath{\mathrm{AF}}\xspace}
\providecommand{\EG}{\ensuremath{\mathrm{EG}}\xspace}
\providecommand{\AG}{\ensuremath{\mathrm{AG}}\xspace}
\providecommand{\OPEN}{\ensuremath{\checkmark}\xspace}
\providecommand{\CLOSED}{\ensuremath{\times}\xspace}
\providecommand{\UNKNOWN}{\ensuremath{\mathsf{Unknown}}\xspace}
\providecommand{\file}[1]{\ensuremath{\texttt{#1}}\xspace}
\providecommand{\func}[1]{\ensuremath{\texttt{#1()}}\xspace}
\providecommand{\Zuniv}{\ensuremath{\top}\xspace}
\providecommand{\Sub}[1]{\ensuremath{\mathit{Sub}(#1)}\xspace}
\providecommand{\Zone}{\ensuremath{\mathcal{Z}}\xspace}
\providecommand{\Ctx}{\ensuremath{\mathcal{X}}\xspace}
\providecommand{\TabGraph}{\ensuremath{\mathcal{G}_{\varphi}}\xspace}
\providecommand{\TabNode}{\ensuremath{\mathcal{N}_{\varphi}}\xspace}
\providecommand{\TabEdge}{\ensuremath{\mathcal{E}_{\varphi}}\xspace}
\providecommand{\nxQ}{\ensuremath{\mathsf{X}_{\mathrm{Q}}}\xspace}
\providecommand{\Sat}{\ensuremath{\mathcal{S}}\xspace}
\providecommand{\FormalSmaller}{\ensuremath{\prec}\xspace}
\providecommand{\tup}[5]{\ensuremath{\langle #1, #2, #3, #4, #5\rangle}\xspace}
\providecommand{\psiU}{\ensuremath{\tEU{p}{q}{[1,2]}}\xspace}
\providecommand{\psiR}{\ensuremath{\tAR{\bot}{p}{[0,1]}}\xspace}
\providecommand{\psiG}{\ensuremath{\AR{\bot}{p}}\xspace}

% Tableau markers and recurring derived notation.
\providecommand{\ANDn}{\ensuremath{\otimes}}
\providecommand{\ORn}{\ensuremath{\oplus}}
\providecommand{\DETn}{\ensuremath{\odot}}
\providecommand{\rst}[2]{#1[#2\,{:=}\,0]}
\providecommand{\kJ}{\ensuremath{J}\xspace}
% Recurrence window of a persistent pursuit: I at the first witness, [0,b-a] after.
\providecommand{\kIrec}{\ensuremath{I^{\ast}}\xspace}
\providecommand{\Inv}{\mathsf{Inv}}
\providecommand{\cvx}{\mathsf{cvx}}
\providecommand{\ZltI}{Z^{\prec\kI}}
\providecommand{\ZinI}{Z^{\in\kI}}
\providecommand{\ZgtI}{Z^{\succ\kI}}
\providecommand{\Dopen}{\Delta^{\circ}}
\providecommand{\Dbd}{\Delta^{\partial}}
\providecommand{\dom}{\operatorname{dom}}
\providecommand{\kind}{\mathit{kind}}
\providecommand{\Acc}{\mathit{Acc}}
\providecommand{\Rej}{\mathit{Rej}}
\providecommand{\FL}{\mathsf{FL}}
\providecommand{\Subt}{\mathsf{T}}

% Algorithm-state shorthands.
\providecommand{\val}{\mathit{val}}
\providecommand{\idx}{\mathit{idx}}
\providecommand{\low}{\mathit{low}}
\providecommand{\onstk}{\mathit{on}}
\providecommand{\ch}{\mathit{ch}}
\providecommand{\Stack}{\mathcal{S}}
\providecommand{\cnt}{\mathit{c}}
\providecommand{\Out}{\mathit{Out}}


%%%%%%%%%%%%%%%%%%%%%%%%%%%%%%%%%%%%%%%%%%%%%%%%%%%%%%%%%%%%%%%%%%%%%%%%%%%%%%
%% Algorithm procedure names.
%% Declared after algorithm2e is loaded so the names type consistently in
%% both the algorithm listings and the prose that refers to them.
%%%%%%%%%%%%%%%%%%%%%%%%%%%%%%%%%%%%%%%%%%%%%%%%%%%%%%%%%%%%%%%%%%%%%%%%%%%%%%
\SetKwFunction{Decide}{Decide}
\SetKwFunction{Explore}{Explore}
\SetKwFunction{Eliminate}{Eliminate}
\SetKwFunction{Expand}{Expand}
\SetKwFunction{Settles}{Settles}
\SetKwFunction{Retire}{Retire}

% Rule keys this presentation adds; the rest are declared in preamble.tex.
\DeclareRule{Carry}{\mathrm{Carry}}
\DeclareRule{Absorb}{\mathrm{Abs}}
\DeclareRule{Umiss}{\mathbf{U}\text{-miss}}
\newcommand{\UstopCarryRuleName}{\mathbf{U}^{\ast}\text{-stop}}
\newcommand{\UunfCarryRuleName}{\mathbf{U}^{\ast}}
\DeclareRule{Ustopcarry}{\UstopCarryRuleName}
\DeclareRule{Uunfcarry}{\UunfCarryRuleName}


%% Algorithms 1 and 2 are tall floats; the LaTeX defaults defer anything over
%% 70% of the text height to a float page at the end of the document, which put
%% them eleven pages after the text that refers to them.
%\renewcommand{\topfraction}{0.92}
%\renewcommand{\bottomfraction}{0.85}
%\renewcommand{\textfraction}{0.06}
%\renewcommand{\floatpagefraction}{0.7}
%\setcounter{topnumber}{3}
%\setcounter{totalnumber}{5}
